\documentclass[11pt,a4paper]{article}
\usepackage[margin=1in]{geometry}
\usepackage{amsmath,amssymb,amsthm}
\usepackage{algorithm}
\usepackage{algpseudocode}
\usepackage{booktabs}
\usepackage{hyperref}

\theoremstyle{definition}
\newtheorem{definition}{Definition}[section]
\newtheorem{theorem}{Theorem}[section]
\newtheorem{lemma}[theorem]{Lemma}
\newtheorem{remark}{Remark}[section]

\title{\textbf{Rigorous Mathematical Specification, Forward-Backward Variance Preservation, and Derivation of Xavier (Glorot) Initialization}}
\author{Team Scaibu \\ \small Email: \texttt{jaydeep.9664920749@gmail.com} \\ \small Architecture and Core Engine Team}
\date{October 2026}

\begin{document}
\maketitle

\begin{abstract}
This document presents the complete theoretical derivation and algorithmic specification of Xavier (Glorot) Initialization. In deep feedforward networks with symmetric linear or sigmoid/tanh activations, random weight initialization with inappropriate variance causes activations to vanish exponentially to zero or explode to infinity across depth. Glorot initialization derives the exact harmonic variance $\sigma^2 = \gamma^2 \frac{2}{n_{\text{in}} + n_{\text{out}}}$ necessary to equalize signal variance during both the forward activation pass and backward gradient backpropagation. We derive the analytical variance balance equations, present the complete algorithmic pseudo-code, step through a numerical calculation, and detail sampling implementations.
\end{abstract}

\section{Notation and Mathematical Preliminaries}

Consider layer $l$ computing linear transformation $y = W x$, where $W \in \mathbb{R}^{n_{\text{out}} \times n_{\text{in}}}$, $x \in \mathbb{R}^{n_{\text{in}}}$, and $y \in \mathbb{R}^{n_{\text{out}}}$.

\begin{itemize}
    \item $n_{\text{in}} \in \mathbb{N}_{\ge 1}$: Input connection fan-in ($n_{\text{in}} = \text{dim}(x)$).
    \item $n_{\text{out}} \in \mathbb{N}_{\ge 1}$: Output connection fan-out ($n_{\text{out}} = \text{dim}(y)$).
    \item $\gamma > 0$: Activation gain factor (e.g., $\gamma = 1.0$ for linear/identity, $\gamma \approx 5/3$ for tanh).
    \item $\sigma^2 \in \mathbb{R}_{> 0}$: Theoretical target variance for weight entries $W_{ij}$.
    \item $\sigma = \sqrt{\sigma^2}$: Standard deviation for Gaussian initialization $W_{ij} \sim \mathcal{N}(0, \sigma^2)$.
    \item $a = \sqrt{3 \sigma^2}$: Half-width bound for uniform initialization $W_{ij} \sim \mathcal{U}(-a, a)$.
\end{itemize}

\begin{table}[h]
\centering
\caption{Summary of Mathematical Symbols for Xavier / Glorot Initialization}
\begin{tabular}{lll}
\toprule
\textbf{Symbol} & \textbf{Type / Domain} & \textbf{Description} \\
\midrule
$n_{\text{in}}$ & $\mathbb{N}_{\ge 1}$ & Fan-in dimension \\
$n_{\text{out}}$ & $\mathbb{N}_{\ge 1}$ & Fan-out dimension \\
$\gamma$ & $\mathbb{R}_{> 0}$ & Activation function gain \\
$\sigma^2$ & $\mathbb{R}_{> 0}$ & Target variance $\gamma^2 \frac{2}{n_{\text{in}} + n_{\text{out}}}$ \\
$\sigma$ & $\mathbb{R}_{> 0}$ & Normal standard deviation \\
$a$ & $\mathbb{R}_{> 0}$ & Uniform half-width bound $\sqrt{3 \sigma^2}$ \\
\bottomrule
\end{tabular}
\end{table}

\section{Problem Definition and Core Concepts}

\subsection{Intuition and Motivation}
When signals propagate forward through $L$ consecutive layers $y = W_L \dots W_1 x$, the variance of the activation at layer $L$ scales as $\operatorname{Var}(y_L) \approx \operatorname{Var}(x) \prod_{l=1}^L (n_{\text{in}}^{(l)} \operatorname{Var}(W^{(l)}))$. If $n_{\text{in}} \operatorname{Var}(W) > 1$, activations explode exponentially; if $< 1$, activations vanish to zero. Simultaneously, during backward backpropagation, gradient variance scales as $\operatorname{Var}(\nabla_x \mathcal{L}) \approx \operatorname{Var}(\nabla_y \mathcal{L}) \prod_{l=1}^L (n_{\text{out}}^{(l)} \operatorname{Var}(W^{(l)}))$. Glorot initialization reconciles both forward ($\operatorname{Var}(W) = 1/n_{\text{in}}$) and backward ($\operatorname{Var}(W) = 1/n_{\text{out}}$) constraints by setting the variance to the harmonic average: $\operatorname{Var}(W) = \frac{2}{n_{\text{in}} + n_{\text{out}}}$.

\subsection{Formal Mathematical Definition}
\begin{definition}[Glorot Normal and Uniform Parameterizations]
Given fan-in $n_{\text{in}}$, fan-out $n_{\text{out}}$, and gain $\gamma > 0$:
\begin{align}
\sigma^2 &= \gamma^2 \frac{2}{n_{\text{in}} + n_{\text{out}}} \label{eq:glorot_var} \\
\sigma &= \sqrt{\sigma^2} = \gamma \sqrt{\frac{2}{n_{\text{in}} + n_{\text{out}}}} \label{eq:glorot_normal} \\
a &= \sqrt{3 \sigma^2} = \gamma \sqrt{\frac{6}{n_{\text{in}} + n_{\text{out}}}} \label{eq:glorot_uniform}
\end{align}
Weights are sampled i.i.d. according to:
\begin{itemize}
    \item \textbf{Glorot Normal}: $W_{ij} \sim \mathcal{N}\left(0, \sigma^2\right)$
    \item \textbf{Glorot Uniform}: $W_{ij} \sim \mathcal{U}\left(-a, a\right)$
\end{itemize}
\end{definition}

\section{Algorithmic Specification}

The computational pipeline implemented in \texttt{NnAlgoXavierGlorotInit} is shown below.

\begin{algorithm}[H]
\caption{Xavier / Glorot Initialization Statistics Computation}
\begin{algorithmic}[1]
\Require Fan-in count $n_{\text{in}} \ge 1$, fan-out count $n_{\text{out}} \ge 1$
\Require Distribution family $\in \{\text{uniform}, \text{normal}\}$, gain $\gamma > 0$
\Ensure Standard deviation $\sigma$, uniform bound $a$, variance $\sigma^2$
\If{$n_{\text{in}} < 1$ \textbf{or} $n_{\text{out}} < 1$ \textbf{or} $\gamma \le 0$}
    \State \textbf{throw} PreconditionFailureException(``Invalid fan dimensions or non-positive gain'')
\EndIf
\State $\text{var} \gets (\gamma^2) \cdot \left(\frac{2.0}{\text{float}(n_{\text{in}} + n_{\text{out}})}\right)$
\State $\text{std} \gets \sqrt{\text{var}}$
\State $\text{bound} \gets \sqrt{3.0 \cdot \text{var}}$
\State \Return $\{\text{std\_dev}: \text{std}, \text{uniform\_bound}: \text{bound}, \text{variance}: \text{var}\}$
\end{algorithmic}
\end{algorithm}

\section{Theoretical Guarantees and Variance Preservation}

\begin{theorem}[Forward and Backward Variance Balance]
Assume input features $x_i$ and weights $W_{ij}$ are mutually independent zero-mean random variables with $\operatorname{Var}(x_i) = \operatorname{Var}(x)$ and $\operatorname{Var}(W_{ij}) = \sigma^2$. For a linear transformation $y_i = \sum_{j=1}^{n_{\text{in}}} W_{ij} x_j$:
\begin{equation}
\operatorname{Var}(y_i) = n_{\text{in}} \sigma^2 \operatorname{Var}(x)
\end{equation}
and for incoming gradient $\frac{\partial \mathcal{L}}{\partial x_j} = \sum_{i=1}^{n_{\text{out}}} W_{ij} \frac{\partial \mathcal{L}}{\partial y_i}$:
\begin{equation}
\operatorname{Var}\left(\frac{\partial \mathcal{L}}{\partial x_j}\right) = n_{\text{out}} \sigma^2 \operatorname{Var}\left(\frac{\partial \mathcal{L}}{\partial y}\right)
\end{equation}
Setting $\sigma^2 = \frac{2}{n_{\text{in}} + n_{\text{out}}}$ enforces $\frac{1}{2}\left(n_{\text{in}} \sigma^2 + n_{\text{out}} \sigma^2\right) = 1.0$, preserving average signal energy.
\end{theorem}

\begin{proof}
By independence of $W_{ij}$ and $x_j$, since $\mathbb{E}[W_{ij}] = 0$:
\begin{align}
\operatorname{Var}(y_i) &= \operatorname{Var}\left(\sum_{j=1}^{n_{\text{in}}} W_{ij} x_j\right) = \sum_{j=1}^{n_{\text{in}}} \operatorname{Var}(W_{ij} x_j) \\
&= \sum_{j=1}^{n_{\text{in}}} \left( \mathbb{E}[W_{ij}^2] \mathbb{E}[x_j^2] - (\mathbb{E}[W_{ij}]\mathbb{E}[x_j])^2 \right) = n_{\text{in}} \sigma^2 \operatorname{Var}(x)
\end{align}
The backward derivation follows identically by swapping the summation index over $n_{\text{out}}$.
\end{proof}

\subsection{Limitations}
\begin{itemize}
    \item \textbf{ReLU Non-Linearity}: Glorot initialization assumes symmetric activations ($f'(0) \approx 1$). For rectified linear units (ReLU), half the activations are zeroed, causing variance to halve at each layer (He/Kaiming initialization must be used instead).
\end{itemize}

\section{Complexity Analysis}

\subsection{Time and Space Complexity}
\begin{itemize}
    \item \textbf{Time}: $\mathcal{O}(1)$ to compute statistics; $\mathcal{O}(n_{\text{in}} \times n_{\text{out}})$ to sample tensor weights.
    \item \textbf{Space}: $\mathcal{O}(1)$ auxiliary memory.
\end{itemize}

\section{Worked Numerical Example}

Consider a linear layer with $n_{\text{in}} = 256$, $n_{\text{out}} = 512$, and $\gamma = 1.0$.

\begin{enumerate}
    \item \textbf{Denominator}:
    \begin{equation}
    n_{\text{in}} + n_{\text{out}} = 256 + 512 = 768
    \end{equation}
    \item \textbf{Target Variance}:
    \begin{equation}
    \sigma^2 = 1.0^2 \times \frac{2.0}{768} = \frac{1}{384} \approx 0.002604167
    \end{equation}
    \item \textbf{Normal Standard Deviation}:
    \begin{equation}
    \sigma = \sqrt{0.002604167} \approx 0.05103104
    \end{equation}
    \item \textbf{Uniform Half-Width Bound}:
    \begin{equation}
    a = \sqrt{3 \times 0.002604167} = \sqrt{0.0078125} \approx 0.08838835
    \end{equation}
\end{enumerate}

\section{Numerical Considerations and Edge Cases}

\begin{itemize}
    \item \textbf{Uniform Variance Identity}: For $\mathcal{U}(-a, a)$, variance is $\frac{(2a)^2}{12} = \frac{a^2}{3}$. Thus setting $a = \sqrt{3 \sigma^2}$ guarantees exact variance equivalence with $\mathcal{N}(0, \sigma^2)$.
    \item \textbf{Conv2D Tensor Extensions}: For 2D convolutions with kernel shape $(C_{\text{out}}, C_{\text{in}}, K_h, K_w)$, $\text{fan\_in} = C_{\text{in}} \cdot K_h \cdot K_w$ and $\text{fan\_out} = C_{\text{out}} \cdot K_h \cdot K_w$.
\end{itemize}

\begin{thebibliography}{9}
\bibitem{glorot2010} X.~Glorot and Y.~Bengio, ``Understanding the difficulty of training deep feedforward neural networks,'' in \emph{International Conference on Artificial Intelligence and Statistics (AISTATS)}, pp.~249--256, 2010.
\bibitem{he2015} K.~He, X.~Zhang, S.~Ren, and J.~Sun, ``Delving deep into rectifiers: Surpassing human-level performance on imagenet classification,'' in \emph{IEEE International Conference on Computer Vision (ICCV)}, pp.~1026--1034, 2015.
\end{thebibliography}

\end{document}
